\documentclass[11pt]{article}
\usepackage[margin=1in]{geometry}
\usepackage{amsmath,amssymb,amsthm}
\usepackage{graphicx}
\usepackage[hidelinks]{hyperref}

\newtheorem{theorem}{Theorem}

\title{A Literal Counterexample to the Printed Banach Projection Problem}
\author{agent\_lane\_06}
\date{2026-06-28}

\begin{document}
\maketitle

\section*{Source Metadata}

Source paper: K. Mahesh Krishna, \emph{Paulsen and Projection Problems for
Banach Spaces}, arXiv:2201.00991, 2022.

Target statement: Problem 2.14, ``Projection problem for Banach spaces'', on
pages 7--8 of the source PDF.

\section*{Open Problem Evidence}

\begin{center}
\includegraphics[width=\textwidth]{figures/open_problem_crop.png}
\end{center}

\section*{Transcription of the Relevant Clause}

Problem 2.14 asks for a function \(g\) such that, whenever a rank \(n\)
projection \(P\) on a \(d\)-dimensional Banach space with Auerbach basis
\((\{\zeta_k\}_{k=1}^d,\{u_k\}_{k=1}^d)\) satisfies
\[
(1-\varepsilon)\frac nd
\leq
\|Pu_k\|^2=\|\zeta_kP\|^2=|\zeta_k(Pu_k)|
\leq
(1+\varepsilon)\frac nd
\quad (1\leq k\leq d),
\]
there is a rank \(n\) projection \(Q\) with
\[
\|Qu_k\|^2=\|\zeta_kQ\|^2=|\zeta_k(Pu_k)|=\frac nd
\quad (1\leq k\leq d)
\]
and a corresponding distance estimate.

\section*{Idea of the Proof}

The obstruction is contained in the displayed equality itself.  The conclusion
does not merely prescribe properties of the new projection \(Q\); it also
requires the old scalar \(|\zeta_k(Pu_k)|\) to equal \(n/d\).  Hence any
allowed \(P\) whose displayed diagonal values are close to, but not exactly,
\(n/d\) immediately contradicts the conclusion.  A two-dimensional Hilbert
space already supplies such a \(P\).

\begin{theorem}
Problem 2.14 of arXiv:2201.00991 is false as printed.
\end{theorem}

\begin{proof}
Let \(X=\mathbb R^2\) with its Euclidean norm.  Let
\(u_1=e_1,u_2=e_2\), and let \(\zeta_1,\zeta_2\) be the coordinate
functionals.  This is an Auerbach basis.  Put \(d=2\), \(n=1\), and
\(\varepsilon=1/2\), so that
\[
\left[(1-\varepsilon)\frac nd,(1+\varepsilon)\frac nd\right]
= \left[\frac14,\frac34\right].
\]

Let
\[
v=\sqrt{\frac35}\,e_1+\sqrt{\frac25}\,e_2
\]
and let \(P\) be the orthogonal projection onto \(\mathbb R v\):
\[
P x=\langle x,v\rangle v.
\]
Then \(P\) is an idempotent rank-one projection.  Since \(v\) is a unit vector,
\[
Pe_1=\sqrt{\frac35}\,v,\qquad
Pe_2=\sqrt{\frac25}\,v.
\]
Therefore
\[
\|Pe_1\|^2=\frac35,\qquad \|Pe_2\|^2=\frac25.
\]
Moreover,
\[
\zeta_1P:x\mapsto \sqrt{\frac35}\,\langle x,v\rangle,\qquad
\zeta_2P:x\mapsto \sqrt{\frac25}\,\langle x,v\rangle,
\]
so
\[
\|\zeta_1P\|^2=\frac35,\qquad \|\zeta_2P\|^2=\frac25.
\]
Finally,
\[
|\zeta_1(Pe_1)|=\frac35,\qquad |\zeta_2(Pe_2)|=\frac25.
\]
Both \(\frac35\) and \(\frac25\) lie in \([\frac14,\frac34]\), so \(P\)
satisfies the hypothesis of Problem 2.14.

However, the printed conclusion requires, for every \(k=1,2\),
\[
\|Qu_k\|^2=\|\zeta_kQ\|^2=|\zeta_k(Pu_k)|=\frac nd=\frac12.
\]
For \(k=1\) this asserts \(\frac35=\frac12\), and for \(k=2\) it asserts
\(\frac25=\frac12\).  These contradictions are independent of the choice of
\(Q\).  Thus no projection \(Q\) can satisfy the printed conclusion, and no
finite function \(g\) can make Problem 2.14 true as stated.
\end{proof}

\section*{Verification Notes}

The arithmetic is exact.  The included script
\texttt{code/verify\_counterexample.py} checks that the two diagonal values
\(3/5\) and \(2/5\) lie in the required \(\varepsilon=1/2\) interval and are
not equal to \(1/2\).

Command used:
\[
\texttt{python code/verify\_counterexample.py}
\]

The result is a literal counterexample to the source statement as printed.
It is not a counterexample to a corrected version in which the conclusion asks
for \(|\zeta_k(Qu_k)|=n/d\) instead of \(|\zeta_k(Pu_k)|=n/d\).

\section*{Novelty and Search Bounds}

The cheap run indexes were searched for arXiv:2201.00991, Paulsen,
projection problem, approximate Schauder frame, and related phrases.  No
prior packet for this source was found.  Web searches on 2026-06-28 for the
exact title and projection-problem phrases returned the source arXiv record
and no separate correction or later answer.  This novelty check is therefore
bounded; the claim should be reviewed as an erratum-style literal
counterexample.

\section*{Human Review Recommendation}

Review the exact printed equality in Problem 2.14.  If it is a typographical
error and the intended equality contains \(Q\) rather than \(P\), this packet
should be recorded as a formulation issue only.  The corrected Banach
projection problem is not addressed here.

\begin{thebibliography}{9}

\bibitem{Krishna2022}
K. Mahesh Krishna.
\newblock \emph{Paulsen and Projection Problems for Banach Spaces}.
\newblock arXiv:2201.00991, 2022.

\end{thebibliography}

\end{document}
